\section{OS/Wightman theorem packet and non-triviality}
\label{app:axioms_package}

This appendix is Packet~9 in the public theorem docket of Section~\ref{sec:roadmap}.
It is invoked directly in Theorem~\ref{thm:main}(C)--(D) and in Appendix~\ref{app:global_form_resolution}.
The constructive inputs remain in Section~\ref{sec:sharp_local} (Lane~A) and Appendix~\ref{app:route1} (Route~1/QE3),
but the point of the present appendix is that a reader can audit the local-field axiomatic output without reverse-engineering the proof spine.
For the shortest referee route through the material collected here, see Appendix~\ref{app:referee_audit}, Extract~C in \S\ref{subsec:hard_module_extracts}, and Stages~6--7 of Appendix~\ref{app:executive_proof_sketch}.
We work throughout with the full local gauge-invariant sharp-local state
\[
\omega^{(u_{\mathrm{ren}})}\in \mathfrak A_{\mathrm{loc}}^*
\]
furnished by Corollary~\ref{cor:LaneA_full_scope} for a fixed weak-window target value $u_{\mathrm{ren}}\in(0,u_\ast]$.
Write $\mathcal P_{+,\mathrm{sharp}}$ for the positive-time sharp-local polynomial algebra from Theorem~\ref{thm:LaneA_bounded_base_cyclic}.
The provenance ledger in \S\ref{subsec:packet9_provenance_ledger} is the front-end audit map for this appendix: it separates imported standard OS/Wightman background from earlier internal Euclidean inputs and from the theorem-level endpoints first packaged here.

\subsection{Packet~9 provenance ledger}
\label{subsec:packet9_provenance_ledger}

For this appendix the provenance bins are used in the following strict sense.
\emph{Imported standard background family} means a standard OS/Wightman reconstruction, analytic-continuation, or uniqueness theorem used as a black-box input;
\emph{earlier internal theorem} means a theorem already proved before Appendix~\ref{app:axioms_package};
and \emph{new appendix endpoint/reformulation} means a theorem-level output first packaged here from those earlier ingredients.

\paragraph{Packet~9 provenance ledger.}
The provenance bins used in this appendix are as follows.
\begin{description}[leftmargin=2.2em,style=nextline]
\item[Theorem~\ref{thm:OS_axioms_package} (new endpoint):] consumes Theorem~\ref{thm:LaneA_closure}(b), Corollary~\ref{cor:LaneA_full_scope}, Theorem~\ref{thm:LaneA_temperedness_verification}, and Theorem~\ref{thm:LaneA_bounded_base_cyclic}, together with the standard OS reconstruction / uniqueness theorem family for item~(v). It packages the full Euclidean OS dossier on the full sharp-local algebra.
\item[Corollary~\ref{cor:OS_cluster_vacuum} (new reformulation):] consumes Theorem~\ref{thm:OS_axioms_package} and Corollary~\ref{cor:route1_QE3_gap}. It converts QE3 into positive-time Euclidean clustering and vacuum uniqueness.
\item[Standard OS reconstruction / uniqueness theorems~\cite{OS1,OS2} (imported background):] used only in Theorem~\ref{thm:OS_axioms_package}, item~(v), and in Theorem~\ref{thm:OS_to_Wightman}.
\item[Theorem~\ref{thm:OS_to_Wightman} (new endpoint):] consumes Theorem~\ref{thm:OS_axioms_package} together with the imported OS reconstruction / analytic-continuation / uniqueness family. It reconstructs the Wightman fields, spectrum condition, and Haag--Kastler net.
\item[Corollary~\ref{cor:field_correspondence} (new endpoint):] consumes Definition~\ref{def:local_operator_spaces}, Theorem~\ref{thm:LaneA_associated_graded_identity}, Corollary~\ref{cor:LaneA_full_scope}, and Theorem~\ref{thm:OS_to_Wightman}. It identifies the reconstructed local fields with the Lane~A renormalized gauge-invariant local differential polynomials.
\item[Corollary~\ref{cor:Minkowski_gap_from_OS} (new endpoint):] consumes Theorem~\ref{thm:OS_to_Wightman} and Corollary~\ref{cor:route1_QE3_gap}. It transfers the OS spectral gap to the Minkowski Hamiltonian.
\item[Theorem~\ref{thm:nontriviality_package} (new endpoint):] consumes Theorem~\ref{thm:LaneA_nontriviality_witness}, Corollary~\ref{cor:field_correspondence}, Corollary~\ref{cor:route1_QE3_gap}, and Corollary~\ref{cor:Minkowski_gap_from_OS}. It rules out the zero/identity theory.
\end{description}


\subsection{Euclidean OS dossier}

\noindent\emph{Direct inputs.} The internal theorem-level input is Theorem~\ref{thm:LaneA_closure}(b), Corollary~\ref{cor:LaneA_full_scope}, Theorem~\ref{thm:LaneA_temperedness_verification}, and Theorem~\ref{thm:LaneA_bounded_base_cyclic}; only the standard OS reconstruction / uniqueness theorem family is imported here, and it enters solely in item~(v).\par

\begin{theorem}[Euclidean OS axioms for the full sharp-local theory]\label{thm:OS_axioms_package}
Let $\omega^{(u_{\mathrm{ren}})}$ be the full local gauge-invariant sharp-local state of Corollary~\ref{cor:LaneA_full_scope}.
Then the Schwinger distributions generated by the renormalized local fields of Lane~A satisfy the following properties on the full sharp-local algebra $\mathfrak A_{\mathrm{loc}}$:
\begin{enumerate}[label=(\roman*),leftmargin=2.3em]
\item Euclidean invariance under translations and rotations.
\item Bosonic symmetry under permutation of insertions.
\item Osterwalder--Schrader reflection positivity on the full positive-time sharp-local algebra.
\item Regularity/temperedness of the Schwinger distributions.
\item OS reconstruction on the full positive-time sharp-local algebra, yielding a Hilbert space
\[
(\mathcal H_{\mathrm{loc}},\Omega_{\mathrm{loc}},H_{\mathrm{loc}}),
\]
with $\Omega_{\mathrm{loc}}=[1]$ cyclic for the full sharp-local positive-time algebra.
Moreover the bounded positive-time base algebra $\mathcal A_+$ embeds densely in $\mathcal H_{\mathrm{loc}}$.
\end{enumerate}
\end{theorem}

\begin{proof}
Items~(i)--(iii) come from the finite-cap sharp-local OS structure theorem, Theorem~\ref{thm:LaneA_closure}(b), together with the inductive-union passage Corollary~\ref{cor:LaneA_full_scope}. Concretely, Euclidean invariance, bosonic symmetry, and reflection positivity are proved on each fixed engineering cap and then transported compatibly to the full sharp-local algebra because every test polynomial uses only finitely many generators. The finite-cap operator/domain package used there is isolated in Lemma~\ref{lem:finite_cap_common_core} and supplies the common invariant core on which the sharp-local field polynomials are realized.

Item~(iv) is Theorem~\ref{thm:LaneA_temperedness_verification}, which verifies the regularity and temperedness bounds for the full sharp-local Schwinger family. Once items~(i)--(iv) are in place, the standard OS reconstruction theorem applies on the full positive-time sharp-local algebra and yields the Hilbert space $(\mathcal H_{\mathrm{loc}},\Omega_{\mathrm{loc}},H_{\mathrm{loc}})$. The additional density statement for $\mathcal A_+$ is then exactly Theorem~\ref{thm:LaneA_bounded_base_cyclic}, which identifies the bounded base algebra as dense inside that reconstructed Hilbert space.
\end{proof}

\dependencyledger
  {Theorem~\ref{thm:LaneA_closure}(b), Corollary~\ref{cor:LaneA_full_scope}, Theorem~\ref{thm:LaneA_temperedness_verification}, and Theorem~\ref{thm:LaneA_bounded_base_cyclic}}
  {the standard OS reconstruction / uniqueness theorem family, used only for item~(v) after the Euclidean dossier has been verified internally}
  {the full sharp-local Schwinger family satisfies the Euclidean OS dossier on the full sharp-local algebra and yields the reconstructed OS Hilbert space with dense bounded base algebra}
  {Corollary~\ref{cor:OS_cluster_vacuum} and Theorem~\ref{thm:OS_to_Wightman}}

\noindent\emph{Direct inputs.} This corollary consumes only Theorem~\ref{thm:OS_axioms_package} and the Route~1 vacuum-gap theorem Corollary~\ref{cor:route1_QE3_gap}.\par

\begin{corollary}[Positive-time Euclidean clustering and vacuum uniqueness]\label{cor:OS_cluster_vacuum}
In the setting of Theorem~\ref{thm:OS_axioms_package}, let
\[
F^0:=F-\omega^{(u_{\mathrm{ren}})}(F),\qquad G^0:=G-\omega^{(u_{\mathrm{ren}})}(G)
\]
for $F,G\in\mathcal P_{+,\mathrm{sharp}}$.
Let $m_*>0$ be the Route~1 gap constant from Corollary~\ref{cor:route1_QE3_gap}.
Then for every $s\ge 0$,
\[
\big|\omega^{(u_{\mathrm{ren}})}(\Theta F^0\,\tau_s G^0)\big|
\le
\|[F^0]\|_{\mathcal H_{\mathrm{loc}}}\,\|[G^0]\|_{\mathcal H_{\mathrm{loc}}}\,e^{-m_* s}.
\]
In particular,
\[
\ker(H_{\mathrm{loc}})=\mathbb C\Omega_{\mathrm{loc}},
\]
so the Euclidean/OS vacuum is unique up to phase.
\end{corollary}

\begin{proof}
Corollary~\ref{cor:route1_QE3_gap} gives
\[
\mathrm{Spec}(H_{\mathrm{loc}})\subset \{0\}\cup[m_*,\infty),
\qquad
\ker(H_{\mathrm{loc}})=\mathbb C\Omega_{\mathrm{loc}}.
\]
For positive-time sharp-local polynomials, OS reconstruction gives the semigroup identity
\[
\omega^{(u_{\mathrm{ren}})}(\Theta F^0\,\tau_s G^0)
=\langle [F^0],e^{-sH_{\mathrm{loc}}}[G^0]\rangle.
\]
Since $[F^0]$ and $[G^0]$ lie in $\Omega_{\mathrm{loc}}^\perp$, the spectral theorem yields
\[
\|e^{-sH_{\mathrm{loc}}}\!\restriction_{\Omega_{\mathrm{loc}}^\perp}\|\le e^{-m_* s},
\]
which gives the stated estimate.
The vacuum-uniqueness statement is the second clause of Corollary~\ref{cor:route1_QE3_gap}.
\end{proof}

\subsection{Wightman/Haag--Kastler reconstruction and field correspondence}

\noindent\emph{Direct inputs.} The new reconstruction endpoint uses Theorem~\ref{thm:OS_axioms_package} as its internal Euclidean dossier and then invokes only the standard Osterwalder--Schrader analytic-continuation / uniqueness theorem family as imported background.\par

\begin{theorem}[OS/Wightman theorem package for the local gauge-invariant sharp-local theory]\label{thm:OS_to_Wightman}
Let $\omega^{(u_{\mathrm{ren}})}$ be as above.
Then the Euclidean Schwinger functions of the full sharp-local theory admit a Minkowski reconstruction with the following properties:
\begin{enumerate}[label=(\roman*),leftmargin=2.3em]
\item there exists a Hilbert space $\mathcal H_{\mathrm{loc}}$, a vacuum vector $\Omega_{\mathrm{loc}}$, and a strongly continuous unitary representation $U$ of the Poincar\'e group whose joint spectrum lies in the closed forward light cone;
\item for every renormalized local gauge-invariant field generated by Lane~A there is a corresponding operator-valued tempered distribution on Minkowski space whose Wightman distributions are the analytic continuations of the Euclidean Schwinger distributions;
\item these fields satisfy locality (spacelike commutativity) and generate the vacuum Haag--Kastler net of the local gauge-invariant sharp-local theory;
\item the reconstruction is unique up to unitary equivalence once the Euclidean Schwinger functions are fixed.
\end{enumerate}
\end{theorem}

\begin{proof}
Theorem~\ref{thm:OS_axioms_package} supplies the full Euclidean dossier on the sharp-local algebra: Euclidean invariance, bosonic symmetry, reflection positivity, temperedness, and the OS-positive positive-time algebra on which reconstruction is to be performed. The standard Osterwalder--Schrader reconstruction and analytic-continuation theorem family therefore applies directly; see~\cite{OS1,OS2}. It yields the reconstructed Wightman fields, the strongly continuous Poincar\'e representation with forward-light-cone spectrum, and spacelike locality after continuation to Minkowski signature.

The Haag--Kastler net is then the vacuum local net generated by the smeared reconstructed fields, and uniqueness up to unitary equivalence is the standard uniqueness clause of the same OS/Wightman reconstruction family once the Euclidean Schwinger functions are fixed. No additional nonstandard reconstruction step is inserted here.
\end{proof}

\dependencyledger
  {Theorem~\ref{thm:OS_axioms_package}}
  {the standard Osterwalder--Schrader analytic-continuation / uniqueness theorem family}
  {the Euclidean sharp-local Schwinger family reconstructs uniquely to the Wightman fields and Haag--Kastler net of the local gauge-invariant sharp-local theory}
  {Corollary~\ref{cor:field_correspondence}, Corollary~\ref{cor:Minkowski_gap_from_OS}, and Theorem~\ref{thm:global_form_local_endpoint}}

\noindent\emph{Direct inputs.} The field-correspondence statement consumes Definition~\ref{def:local_operator_spaces}, Theorem~\ref{thm:LaneA_associated_graded_identity}, Corollary~\ref{cor:LaneA_full_scope}, and the reconstruction theorem Theorem~\ref{thm:OS_to_Wightman}.\par

\begin{corollary}[Field correspondence for gauge-invariant local differential polynomials]\label{cor:field_correspondence}
Let $[P]$ be any exact-dimension quotient class from Definition~\ref{def:local_operator_spaces}, and let
\[
\mathcal O^{\mathrm{ren}}_{[P]}(f;\mu)
\]
denote the corresponding renormalized sharp-local field constructed by Lane~A.
Then in the Wightman reconstruction of Theorem~\ref{thm:OS_to_Wightman} there exists a unique operator-valued tempered distribution
\[
\Phi_{[P]}
\]
whose Euclidean Schwinger distributions are the analytic continuations of the mixed Schwinger functions of $\mathcal O^{\mathrm{ren}}_{[P]}(f;\mu)$.
This correspondence is compatible with finite linear combinations, derivatives, total-derivative/Bianchi/symmetry quotient relations, finite engineering truncations, and the inductive-union passage to the full sharp-local algebra.
Consequently the reconstructed local field content is exactly the family of gauge-invariant local differential polynomials in the curvature and its covariant derivatives represented by the Lane~A sharp-local algebra.
\end{corollary}

\begin{proof}
Definition~\ref{def:local_operator_spaces} and Theorem~\ref{thm:LaneA_associated_graded_identity} identify the exact-dimension quotient classes before any basis is chosen, and Corollary~\ref{cor:LaneA_full_scope} constructs compatible renormalized fields for every finite truncation and hence on the inductive union.
Applying Theorem~\ref{thm:OS_to_Wightman} to the Euclidean Schwinger functions of those renormalized fields gives the corresponding Minkowski operator-valued tempered distributions.
Because the quotient relations and truncation compatibilities are already built into the Euclidean theory before reconstruction, the same relations persist after analytic continuation.
\end{proof}

\noindent\emph{Direct inputs.} This gap-transfer corollary consumes only Theorem~\ref{thm:OS_to_Wightman} and the Route~1 spectral-gap statement Corollary~\ref{cor:route1_QE3_gap}.\par

\begin{corollary}[Minkowski energy gap from the OS spectral gap]\label{cor:Minkowski_gap_from_OS}
Let $P^0$ be the Minkowski Hamiltonian in the Wightman reconstruction of Theorem~\ref{thm:OS_to_Wightman}.
Assume the Route~1 vacuum-gap theorem, i.e. Corollary~\ref{cor:route1_QE3_gap}, so that
\[
\mathrm{Spec}(H_{\mathrm{loc}})\subset \{0\}\cup[m_*,\infty).
\]
Then
\[
\mathrm{Spec}(P^0)\subset \{0\}\cup[m_*,\infty).
\]
In particular, the reconstructed local gauge-invariant sharp-local theory has positive energy, a unique vacuum, and a strictly positive Minkowski mass gap in its vacuum sector.
\end{corollary}

\begin{proof}
Under the Osterwalder--Schrader reconstruction, Euclidean time translations analytically continue to the Minkowski time-translation group.
Thus the Minkowski Hamiltonian $P^0$ coincides with the OS generator $H_{\mathrm{loc}}$ on the reconstructed Hilbert space; see Theorem~\ref{thm:OS_to_Wightman}.
The spectral inclusion therefore transfers directly from Corollary~\ref{cor:route1_QE3_gap}.
\end{proof}


\noindent\emph{Direct inputs.} The non-triviality package consumes Theorem~\ref{thm:LaneA_nontriviality_witness}, Corollary~\ref{cor:field_correspondence}, Corollary~\ref{cor:route1_QE3_gap}, and Corollary~\ref{cor:Minkowski_gap_from_OS}; only the final K\"all\'en--Lehmann reformulation uses standard background language.\par

\begin{theorem}[Non-triviality witness and nonzero spectral weight]\label{thm:nontriviality_package}
Assume the full sharp-local OS/Wightman construction of Theorem~\ref{thm:OS_to_Wightman} together with the Route~1 vacuum-gap theorem
Corollary~\ref{cor:route1_QE3_gap}.
Let $\mathcal E_{\mathrm{ren}}$ be the canonical centered CP-even dimension-$4$ scalar witness from Theorem~\ref{thm:LaneA_nontriviality_witness},
and let $\phi_{\ast,S}$ be the corresponding positive-time reference test function.
Then:
\begin{enumerate}[label=(\roman*),leftmargin=2.3em]
\item
\[
\|\mathcal E_{\mathrm{ren}}(\phi_{\ast,S})\Omega_{\mathrm{loc}}\|^2
=\omega^{(u_{\mathrm{ren}})}\big(\Theta\mathcal E_{\mathrm{ren}}(\phi_{\ast,S})\,\mathcal E_{\mathrm{ren}}(\phi_{\ast,S})\big)=1;
\]
\item the centered Euclidean and Minkowski two-point functions of $\mathcal E_{\mathrm{ren}}$ are therefore nonzero;
\item the spectral measure of the nonzero vector $\mathcal E_{\mathrm{ren}}(\phi_{\ast,S})\Omega_{\mathrm{loc}}$ under both $H_{\mathrm{loc}}$ and the Minkowski Hamiltonian $P^0$ is a nonzero finite positive measure supported in $[m_\ast,\infty)$.
Equivalently, the K\"all\'en--Lehmann measure of the centered two-point function of the witness field is not the zero measure.
\end{enumerate}
In particular, the reconstructed local gauge-invariant sharp-local theory is non-trivial in the load-bearing sense required later in the paper: it contains a local gauge-invariant field that is not a multiple of the identity and has nonzero vacuum-created spectral weight above the gap.
\end{theorem}

\begin{proof}
The norm identity in~(i) is exactly Theorem~\ref{thm:LaneA_nontriviality_witness}, interpreted in the OS Hilbert space.
Because $\mathcal E_{\mathrm{ren}}$ is centered, the vector $\mathcal E_{\mathrm{ren}}(\phi_{\ast,S})\Omega_{\mathrm{loc}}$ is orthogonal to the vacuum.
Hence Corollary~\ref{cor:route1_QE3_gap} and Corollary~\ref{cor:Minkowski_gap_from_OS} place its Euclidean and Minkowski spectral measures inside $[m_\ast,\infty)$.
Since the vector has norm one, those measures are nonzero finite positive measures.
If the centered Euclidean or Minkowski two-point function were identically zero, the corresponding vacuum-created vector would have zero norm, contradicting~(i).
The K\"all\'en--Lehmann reformulation is the standard spectral representation of the centered Wightman two-point function of a local field in the vacuum representation, applied here to the witness field furnished by Corollary~\ref{cor:field_correspondence}.
Thus the theory is non-trivial in the claimed sense.
\end{proof}
